\section{Introduction}
Self-improving agents increasingly convert completed episodes into persistent natural-language memory. This creates an update channel that sits between ordinary evaluation and parameter learning: a terminal label can decide what durable state the agent carries into later tasks. Existing work has studied unreliable judges and reward models~\citep{tan2025judgebench,lambert2024rewardbench}, outcome-conditioned reflection~\citep{shinn2023reflexion,ouyang2026reasoningbank}, learned memory construction~\citep{wang2025memalpha,li2026attrimem}, and reward-dependent memory valuation or reuse~\citep{asadolahi2026inflation,yang2026romerl}. A missing question is more basic: \emph{when a reward-conditioned writer changes persistent state, how much of that state difference actually survives the later memory pipeline and becomes behavior?}

This distinction matters because current evaluation practice often collapses several different objects. A memory representation may change strongly while never being retrieved. A retrieved branch-specific memory may enter context without changing the next action. A changed action may still leave the terminal task outcome unchanged. Conversely, a forced-memory test may reveal that two memories are behaviorally capable of producing different outcomes even though the native retrieval-and-uptake pipeline rarely realizes that leverage. Treating any one of these measurements as a proxy for the others can misdiagnose both the mechanism and the appropriate intervention.

We study this boundary in the released ReasoningBank implementation. Its writer exposes a particularly clean intervention surface: score one selects a success-reflection writer and score zero selects a failure-reflection writer over the same trajectory. We hold the source trajectory fixed and flip only this reward-conditioned writer branch. This isolates an operational write-time state contrast without learning a new construction policy. A disjoint same-mode semantic-rewording control then tests the strongest simple reduction---that the observed state difference is merely ordinary prompt sensitivity.

The upstream effect is robust. Across 20 complete Shopping success/failure writer pairs, the persistent memories diverge. In the frozen eight-trajectory prompt control, reward-mode memory divergence is 0.700 versus 0.595 under stronger semantic-preserving same-mode wording perturbations; the paired excess is 0.105 with exact $p=0.0078$. Thus the reward/reflection branch is a reliable state intervention at write time.

The downstream result is not a monotone continuation of that effect. When paired memories are directly forced into a fixed-evidence future-task packet, they are capable of changing terminal behavior: the fully crossed 4-by-4 confirmatory experiment yields mean absolute success-rate difference $0.15625$ with permutation $p=0.00074$. But this is a \emph{leverage} test, not a native-pipeline transport estimate. Under native Shopping reuse, branch-specific memory is retrieved on 125/172 held-out opportunities, yet the mean absolute terminal success difference is only $0.02083$ ($p=0.4289$), with 34/36 cells exactly zero. The corresponding first-action distribution has TV distance $0.06944$ ($p=0.5801$), and no frozen state changes its modal first action across reward branches. A Reddit replication again changes memory in 4/4 paired writes, but native terminal transport remains sparse and sign-heterogeneous: mean absolute difference $0.125$, $p=0.2253$, with 6/8 cells zero and opposite signs in the two nonzero cells.

These measurements change the scientific object. The paper is not a claim that reward error uniformly contaminates future behavior, nor a claim that memory distance is a behavioral metric. Instead, it identifies a \emph{stage-resolved transport boundary}:
\begin{equation}
\text{reward label} \rightarrow \text{durable write} \rightarrow \text{retrieval exposure} \rightarrow \text{policy uptake} \rightarrow \text{outcome}.
\end{equation}
The treatment is strong at the first stage, demonstrably has downstream leverage when forced, but is sharply attenuated under native exposure and uptake. The mismatch between these measurements is itself the central finding.

\paragraph{Contributions.}
We make three tightly coupled contributions. First, we formulate reward-conditioned memory as a controlled persistent-state intervention by holding the source trajectory fixed and varying only the released success/failure writer branch; a stronger same-mode wording control separates this intervention from generic prompt sensitivity. Second, we introduce a stage-resolved evaluation that distinguishes \emph{state divergence}, \emph{forced behavioral leverage}, \emph{native retrieval exposure}, \emph{policy uptake}, and \emph{terminal transport}; the experiments show that these quantities can differ by an order of magnitude and should not be used interchangeably. Third, we replicate the write/transport separation across Shopping and Reddit and derive a practical evaluation rule: persistent-memory systems should localize where an update's effect attenuates before assigning it behavioral authority or designing a repair.

\paragraph{Closest-work boundary.}
Prior work already establishes that rewards can shape memory construction, valuation, or reuse. We surrender those broad claims. Our residual contribution is measurement and identification: under a fixed writer policy and same-trajectory intervention, durable branch-specific state is easy to create, while native behavioral transport is sparse and domain/task dependent. This is distinct from learning a better memory policy or proposing another memory gate. Indeed, our attempted evidence-gated method extension stopped before behavioral evaluation because its semantic-validity signal could not be independently qualified; we therefore keep the present paper centered on the measured transport boundary rather than manufacturing a solution after seeing the results.
